\documentclass[letterpaper, 12pt, titlepage]{article} % Tipo de documento

% ================================================================
% PAQUETES BÁSICOS Y CONFIGURACIÓN ESENCIAL
% ================================================================
\usepackage[utf8]{inputenc}        % Codificación de caracteres UTF-8
\usepackage[spanish]{babel}        % Localización en español (guiones, títulos)
\usepackage[margin=1in]{geometry}  % Configuración de márgenes de página

% ================================================================
% TIPOGRAFÍA Y FORMATO DE TEXTO
% ================================================================
\usepackage{helvet}                % Fuente Sans-Serif (Helvetica-like)
\renewcommand{\familydefault}{\sfdefault} % Establece Helvet como fuente principal
\usepackage{setspace}              % Control de interlineado (1.5, doble espacio)
\usepackage{parskip}               % Espaciado entre párrafos (en lugar de sangría)
\usepackage{ragged2e}              % Mejor manejo de justificación de texto

% ================================================================
% ELEMENTOS GRÁFICOS Y DISEÑO
% ================================================================
\usepackage{graphicx}              % Manejo de imágenes (inserción y escalado)
\usepackage{xcolor}                % Sistema de colores personalizados
\usepackage{pdflscape}             % Páginas en orientación horizontal

% ================================================================
% TABLAS PROFESIONALES
% ================================================================
\usepackage{booktabs}              % Reglas tipográficas para tablas (formato APA)
\usepackage{xltabular}             % Tablas largas que ocupan todo el ancho del texto
% ================================================================
% LISTAS Y ESTRUCTURAS
% ================================================================
\usepackage{enumitem}              % Personalización avanzada de listas
\setlist{nosep, leftmargin=*, wide=0pt} % Listas ultra compactas

% ================================================================
% HIPERVÍNCULOS Y METADATOS PDF
% ================================================================
\usepackage{hyperref}              % Enlaces interactivos y metadatos PDF
\hypersetup{
    colorlinks = true,             % Colores en lugar de cajas
    linkcolor  = blue!70,          % Color para enlaces internos
    urlcolor   = cyan!80,          % Color para URLs
    citecolor  = magenta!70,           % Color para citas
}

% ================================================================
% CONFIGURACIONES FINALES
% ================================================================
\renewcommand{\thesection}{\Roman{section}} % Numeración romana para secciones
\newcommand{\universidad}{Universidad de las Regiones Autónomas de la Costa Caribe Nicaragüense} % Nombre de la universidad
\newcommand{\carrera}{Nombre de la carrera} % Nombre de la asignatura
\newcommand{\asignatura}{Nombre de la asignatura} % Nombre de la asignatura
\newcommand*{\thead}[1]{\multicolumn{1}{>{\centering\arraybackslash}X}{\bfseries #1}} % Encabezado centrado y en negritas en tablas

\begin{document}

\begin{titlepage}
    \centering
    \begin{minipage}{0.25\textwidth}
        \centering
        \includegraphics[width=0.6\textwidth]{images/uraccan_logo.png}
    \end{minipage}%
    \begin{minipage}{0.75\textwidth}
        \centering
        \Large{\textbf{\universidad}}\\
    \end{minipage}

    \vfill
    \large{\textbf{Plan de clase}}\\
    \Huge{\textbf{\asignatura}}\\
    \vfill
\end{titlepage}
\section{Datos Generales}

\begin{xltabular}{\linewidth}{@{}>{\bfseries}X X@{}}
    \toprule
    Número de clase:                           & Sesión XX de XX               \\
    \midrule
    Fecha:                                     & 03 de marzo de 2025           \\
    \midrule
    Carrera:                                   & Nombre de la carrera          \\
    \midrule
    Asignatura:                                & Nombre de la asignatura       \\
    \midrule
    Año académico:                             & XX año                        \\
    \midrule
    Duración:                                  & 4 horas                       \\
    \midrule
    Turno:                                     & Matutino/Vespertino/Nocturno  \\
    \midrule
    Unidad a desarrollar:                      & Unidad X: Nombre de la unidad \\
    \midrule
    Tema(s):                                   & \begin{enumerate}[topsep=0pt, partopsep=0pt, nosep]
        \item Tema 1
        \item Tema 2
        \item Tema 3
        \item Tema 4
    \end{enumerate}           \\
    \midrule
    Subtema(s):                                & \begin{itemize}
        \item Subtema 1.1
        \item Subtema 1.2
        \item Subtema 1.3
        \item Subtema 1.4
    \end{itemize}        \\
    \midrule
    Objetivos:                                 & \begin{itemize}
        \item Objetivo 1
        \item Objetivo 2
        \item Objetivo 3
        \item Objetivo 4
        \item Objetivo 5
    \end{itemize}               \\
    \midrule
    Formas de Organización de Enseñanzas (FOE): & \begin{itemize}
        \item Estrategia metodológica 1
        \item Estrategia metodológica 2
        \item Estrategia metodológica 3
        \item Estrategia metodológica 4
    \end{itemize}                               \\
    \midrule
    Medios de enseñanza:                       & \begin{itemize}
        \item Material o recurso 1
        \item Material o recurso 2
        \item Material o recurso 3
        \item Material o recurso 4
    \end{itemize}                               \\
    \midrule
    Bibliografía:                              & \begin{itemize}
        \item Autor 1, Título 1, Editorial 1, Año 1.
        \item Autor 2, Título 2, Editorial 2, Año 2.
        \item Autor 3, Título 3, Editorial 3, Año 3.
        \item Autor 4, Título 4, Editorial 4, Año 4.
    \end{itemize}                               \\                                         \\
    \bottomrule
\end{xltabular}
\pagebreak
\section{Introducción}

Para iniciar, el docente se presentará brevemente, compartiendo su experiencia y el propósito de la sesión. Posteriormente, se realizará una evaluación diagnóstica para identificar los conocimientos previos de los estudiantes y ajustar el enfoque de la clase según sea necesario.

A continuación, se presentarán los contenidos que se desarrollarán durante la sesión, destacando su relevancia y cómo estos se vinculan al desarrollo académico y profesional de los estudiantes. Asimismo, se expondrán los objetivos específicos de la clase, los cuales guiarán el aprendizaje y las actividades a realizar.

A lo largo de la clase, se fomentará la participación activa y el análisis crítico para enriquecer el aprendizaje colectivo, asegurando que los contenidos contribuyan al crecimiento integral de los estudiantes.

\subsection*{Preámbulo (5 minutos)}
\begin{itemize}
    \item Aseguramiento del orden, puntualidad, limpieza y disciplina en el aula.
    \item Breve bienvenida a los estudiantes.
    \item Toma de asistencia y verificación de la lista de participantes.
\end{itemize}

\subsection*{Introducción (10 minutos)}
\begin{itemize}
    \item Presentación del contenido de la clase.
    \item Presentación de los objetivos de aprendizaje y estructura de la sesión.
    \item Anuncio de actividades y herramientas a utilizar.
    \item Contextualización de la importancia del contenido de la sesión para el desarrollo académico y profesional.
    \item Pautas de participación y colaboración.
\end{itemize}

\subsection*{Evaluación diagnóstica (15 minutos)}
\begin{itemize}
    \item Evaluación inicial para identificar el nivel de conocimientos previos de los estudiantes.
    \item Uso de herramientas digitales (ej: Kahoot, Google Forms).
    \item Discusión grupal para compartir respuestas y contrastar opiniones.
    \item Retroalimentación inmediata para aclarar dudas y reforzar aprendizajes.
\end{itemize}
\section{Desarrollo}

\subsection*{Motivación y contextualización inicial (5 minutos)}
\begin{itemize}
    \item Presentación de un caso real relacionado con el tema para generar discusión inicial.
    \item Pregunta detonante: "¿Cómo resolverían X problema en su ámbito profesional?".
\end{itemize}

\subsection*{Conferencia (60 minutos)}
\begin{itemize}
    \item Explicación de conceptos clave mediante una presentación interactiva, invitando a los estudiantes a compartir ejemplos.
    \item Uso de herramientas digitales (ej: Google Slides, Prezi, Mentimeter).
    \item Inclusión de ejemplos prácticos para ilustrar los conceptos.
    \item Preguntas abiertas para fomentar la participación.
\end{itemize}

\subsection*{Taller práctico (60 minutos)}
\begin{itemize}
    \item Trabajo práctico guiado para aplicar los conocimientos teóricos adquiridos y ejercitar los prácticos en un caso, una simulación o un laboratorio virtual.
    \item Uso de herramientas digitales (ej: Google Jamboard para mapas conceptuales colaborativos)
    \item Verificación formativa (15 minutos):
    \item Ejercicio rápido de autoevaluación (ej: formulario en Mentimeter con preguntas de opción múltiple).
    \item Discusión grupal para aclarar dudas y reforzar aprendizajes.
\end{itemize}
\section{Cierre}

\subsection*{Resumen (15 minutos)}
\begin{itemize}
    \item Síntesis de los puntos clave del contenido de la sesión con participación de los estudiantes.
    \item Debate sobre los temas tratados.
    \item Preguntas y respuestas.
\end{itemize}

\subsection*{Verificación de objetivos (10 minutos)}
\begin{itemize}
    \item Dominio de los conceptos clave: \textit{¿Qué aprendieron hoy?}
    \item Aplicación de los conceptos en situaciones prácticas: \textit{¿Cómo pueden aplicar lo aprendido en su vida académica y profesional?}
    \item Resolución de problemas y toma de decisiones: \textit{¿Cómo pueden resolver problemas similares en el futuro?}
    \item Evaluación de habilidades y competencias adquiridas: \textit{¿Qué habilidades desarrollaron durante la sesión?}
    \item Reflexión sobre la relación entre los contenidos y la realidad: \textit{¿Cómo se relacionan los temas tratados con su entorno?}
    \item Reflexión sobre el aprendizaje: \textit{¿Qué fue lo más relevante de la sesión?}
    \item Retroalimentación sobre la metodología y dinámica de la clase: \textit{¿Qué les pareció la forma en que se desarrolló la sesión?}
\end{itemize}

\subsection*{Asignación de trabajo independiente (5 minutos)}
\begin{itemize}
    \item Lectura obligatoria: capítulo de libro, artículo académico, informe técnico.
    \item Investigación individual: búsqueda de información sobre un tema específico o análisis de caso, presentado en un informe, ensayo, exposición o entrada de blog.
    \item Ejercicios prácticos: aplicación de conocimientos mediante la resolución de ejercicios, problemas o casos prácticos.
    \item Práctica de laboratorio: diseño, aplicación, resolución de problemas e implementación de tecnologías, experimentos o simulaciones.
\end{itemize}

\end{document}